\begin{helper}
For positive integers \(s\) and \(k\), the two-color Ramsey number
\(\noderef{RamseyNumber}(s,k)\) is positive.
\end{helper}
\begin{proof}
By \(\noderef{RamseyNumber}\), the Ramsey number is the least integer \(n\)
for which the Ramsey property \(R(s,k,n)\) holds.  The set of such \(n\) is
nonempty by \(\noderef{RamseyPropertyNonempty}\), so this least integer belongs
to the defining set.

It remains to show that \(R(s,k,0)\) is false.  Consider the unique graph on
the empty vertex set.  A clique of size \(s\) in this graph would be a finite
subset of the empty set with cardinality \(s\), but every subset of the empty
set has cardinality \(0\), contradicting \(s>0\).  Similarly, an independent
set of size \(k\) would be a subset of the empty set with cardinality \(k\),
contradicting \(k>0\).  Thus \(R(s,k,0)\) fails.

If \(\noderef{RamseyNumber}(s,k)=0\), then the least element of the nonempty
defining set would witness \(R(s,k,0)\), contradicting the previous paragraph.
Therefore \(\noderef{RamseyNumber}(s,k)>0\).
\end{proof}
